零样本异常检测（Zero-Shot Anomaly Detection, ZSAD）借助 CLIP 的开放词表对齐能力，无需目标域样本即可检测缺陷，具有很强的实用价值。然而 CLIP 的对比预训练只优化物体级的全局语义对齐，其 patch 特征在频率上偏低通：尽管异常确实会扰动特征的高频分量，但正常的粗糙材质（如布料、编织、网格）本身也具有很强的高频响应。因此，直接从 CLIP 读出的高频信号是\textbf{歧义}的——高频响应大，既可能是异常，也可能只是材质本身。本文提出一个关键观察：异常相关的高频线索其实已经存在于 CLIP 特征中，缺的不是能力，而是一个描述“这张图正常长什么样”的参照。基于此，我们提出一种训练无关的方法：用 Haar 小波从 CLIP patch 特征中读出高频；用这个独立于 CLIP 语义的信号选择可信的正常/异常证据；进而在\textbf{每张测试图上估计一个正常参照}，并据此重算异常图，将异常定义为高频信号相对该逐图参照的偏离。整个流程不训练、不反传、不使用任何辅助标注数据。在五个数据集上的实验表明：直接融合高频而不建参照会使像素级性能崩盘（pixel AUROC 由 $92.4$ 跌至 $82.5$）；使用全局固定参照明显劣于逐图参照；我们的方法能召回约 $18\text{--}22\%$ 被 CLIP 独自漏检的异常，且增益恰好集中在 CLIP 频率盲的纹理与微缺陷类别上。
